\section{Active normal cones: complete support proofs and a structural obstruction}
\label{sec:active-faces}

The previous continuation replaces some normal-ray enumeration queries by
exact positive-Euler linear programmes and certifies their negative answers
by Farkas multipliers.  It also identifies a difficulty with a strong
mandatory-type backdoor: a tetrahedron unused by a positive surface can
retain several allowed types, even when that surface is easy to find
\cite{priorDual}.  An immediate response is to distinguish an allowed
coordinate from one that is actually positive somewhere in the relaxed cone.
This section makes that distinction exact, implements its certificate, and
then proves a limitation of the resulting structural parameter.

The linear-programming ingredients are classical.  Strict complementarity
and maximal support belong to the Goldman--Tucker framework
\cite{GoldmanTucker}; coordinate-face traversal and normal-surface LP search
have substantial precedents \cite{BurtonFaces,BurtonOzlen}.  Our purpose is
to give a source-sized proof interface, an explicit homogeneous construction
compatible with the supplied exact solver, and a precise complexity bound
for the adaptive tree.  Familiar local matching identities then yield a
linear lower bound for its root parameter on genuine normal-coordinate
inputs, including the maintained diagram-exterior construction.

\subsection{One primal--dual pair identifies every active coordinate}

Let $A\in\Q^{m\times N}$ and $c\in\Q^N$.  Put
\begin{equation}\label{eq:active-cone-def}
 C=\{x\in\R^N:x\ge0,\ Ax=0\},\qquad
 S(C)=\{j:\text{some }z\in C\text{ has }z_j>0\}.
\end{equation}
The set $S(C)$ is the \emph{active support of the cone}; it need not equal the
support of a particular basic solution.  Write
$P_c=\{x\in C:c^Tx>0\}$.  If $P_c$ is nonempty, its active support is also
$S(C)$.  Indeed, if $p\in P_c$ and $z\in C$ has $z_j>0$, then
$p+\varepsilon z\in P_c$ for sufficiently small positive $\varepsilon$.
The objective of $z$ may be negative.  Activity on the positive-objective
part therefore cannot be inferred from the sign of an individual direction's
objective.

\begin{lemma}[Complete active-support alternative]\label{lem:active-pair}
Exactly one of the following alternatives holds:
\begin{enumerate}
\item There are rational vectors $x\in\Q^N$ and $y\in\Q^m$ satisfying
\begin{equation}\label{eq:active-primal-dual}
 Ax=0,\quad x\ge0,\quad c^Tx\ge1,\qquad
 s:=A^Ty\ge0,\quad x+s\ge\mathbf1.
\end{equation}
\item There is a rational vector $h\in\Q^m$ such that
\begin{equation}\label{eq:active-negative}
 A^Th\ge c.
\end{equation}
\end{enumerate}
In the first alternative, $x_j\ge1$ exactly when $j\in S(C)$, and
$s_j\ge1$ exactly when $j\notin S(C)$.
\end{lemma}

\begin{proof}
For a pair satisfying \eqref{eq:active-primal-dual},
\[
 x^Ts=x^TA^Ty=(Ax)^Ty=0.
\]
Both vectors are nonnegative.  Hence $x_js_j=0$ in every coordinate, while
$x_j+s_j\ge1$ ensures that exactly one of the two entries is at least one.
For any $z\in C$, the same calculation gives $z^Ts=0$.  Thus $s_j>0$
forces $z_j=0$ throughout $C$, whereas $x_j>0$ explicitly demonstrates
activity.  This proves the support assertion.

Suppose $P_c\ne\varnothing$.  For each $j\in S(C)$ choose a rational
$z^{(j)}\in C$ with $z^{(j)}_j>0$.  Rational witnesses exist because the
relevant linear systems are rational.  Their sum $z$ is positive on all of
$S(C)$.  Starting from a rational $p\in P_c$, choose a sufficiently small
positive rational $\varepsilon$ and then rescale $p+\varepsilon z$ to obtain
$x_j\ge1$ on $S(C)$ and $c^Tx\ge1$.

For every inactive coordinate $j$, the functional $e_j^Tz$ is nonpositive
on $C$; in fact it is identically zero.  The homogeneous Farkas alternative
gives a rational $y^{(j)}$ with $A^Ty^{(j)}\ge e_j$.  Let
$y=\sum_{j\notin S(C)}y^{(j)}$, with the empty sum interpreted as zero.
Its residual is nonnegative everywhere and at least one off $S(C)$.
Orthogonality with the constructed $x$ makes it zero on $S(C)$.  Thus
\eqref{eq:active-primal-dual} holds.

If $P_c=\varnothing$, homogeneous Farkas gives
\eqref{eq:active-negative}.  Conversely, pairing that inequality with any
nonnegative kernel vector gives $c^Tx\le h^TAx=0$, which excludes the first
alternative.
\end{proof}

The proof has a useful verification consequence.  The checker does not need
one infeasibility test per inactive coordinate, and it does not reproduce a
relative-interior algorithm.  It checks two matrix products, signs, and the
coordinatewise coverage inequality in \eqref{eq:active-primal-dual}.  One
multiplier exposes all inactive coordinates simultaneously; one positive
vector proves that no remaining coordinate has been silently omitted.

The claim is about the relaxed cone.  For example,
\[
 A=(1,-1),\qquad c=(1,1),\qquad C=\{(a,a):a\ge0\}
\]
has both coordinates active and positive objective away from zero.  If the
two coordinates are declared incompatible, it has no admissible positive
point.  A complete active-support vector can thus be a particularly
incompatible relaxation vector.

\subsection{A single homogeneous query with a small negative output}

The inherited solver accepts positive-objective queries over a nonnegative
homogeneous kernel.  We can use that interface directly.  Introduce
nonnegative variables
\[
 (x,y^+,y^-,s,u,v,\tau),
\]
where $x,s,u$ have $N$ entries, $y^+,y^-$ have $m$ entries, and $v,\tau$ are
scalars.  Impose
\begin{align}
 Ax&=0, &
 s-A^T(y^+-y^-)&=0,\label{eq:active-lift-a}\\
 x+s-u-\tau\mathbf1&=0, &
 c^Tx-v-\tau&=0.\label{eq:active-lift-b}
\end{align}
There are $3N+2m+2$ nonnegative variables and $m+2N+1$ equations.  The
objective is $\tau$.  A point with $\tau>0$, divided by $\tau$, yields
\eqref{eq:active-primal-dual} with $y=y^+-y^-$.  Conversely, split any
multiplier from that certificate into its positive and negative parts, set
$\tau=1$, and use the nonnegative surpluses for $u$ and $v$.  The lifted
positive query is therefore equivalent to $P_c\ne\varnothing$.

It is also possible to discard the large lifted dual after a negative
answer.  The following projection is useful because the consumer can remain
an ordinary homogeneous-Farkas checker on the original matrix.

\begin{proposition}[Projection of a lifted negative proof]
\label{prop:active-dual-projection}
Let $(\alpha,\beta,\gamma,\delta)$ be a negative Farkas multiplier for the
lifted positive-$\tau$ query, grouped according to the four displayed sets of
equations.  Then $\delta<0$, and $h=\alpha/(-\delta)$ satisfies
$A^Th\ge c$.
\end{proposition}

\begin{proof}
The lifted column inequalities are
\begin{align*}
 A^T\alpha+\gamma+\delta c&\ge0,& A\beta&=0,\\
 \beta+\gamma&\ge0,& \gamma&\le0,\\
 \delta&\le0,&-\mathbf1^T\gamma-\delta&\ge1.
\end{align*}
The equality $A\beta=0$ comes from the two oppositely signed columns for
each free dual variable.  If $\delta=0$, then
$\beta\ge-\gamma\ge0$.  Pairing the first inequality with $\beta$ gives
\[
 0\le\beta^T\gamma\le-\sum_j\gamma_j^2.
\]
Thus $\gamma=0$, which contradicts
$-\mathbf1^T\gamma-\delta\ge1$.  Hence $\delta<0$.  Finally,
\[
 A^T\alpha\ge-\gamma-\delta c\ge-\delta c,
\]
and division by $-\delta$ gives the required original-cone dual.
\end{proof}

The enlarged system has polynomial encoding length in the original rational
input.  Standard rational-LP algorithms consequently decide it in polynomial
bit complexity, with polynomial-bit primal or dual witnesses.  The executable
producer uses the exact Bland-simplex routine inherited from the previous
continuation \cite{priorDual}.  Its arithmetic is exact and its configured
pivot allowance is explicit; this implementation is not asserted to have a
polynomial worst-case pivot count.  The distinction is material: one large
degenerate LP can cost much more than several small support probes.  The
experiments section reports that tradeoff, including unfavorable outcomes.

\subsection{An intrinsic rank for the remaining incompatibilities}

Let $\mathcal G$ be a collection of coordinate conflict groups.  In normal
coordinates these are the three quadrilateral types in each tetrahedron.
A vector is admissible when each group contains at most one positive
coordinate.  The following definition also permits overlapping groups, since
the proof uses only their incompatibility disjunctions.

\begin{definition}[Active ambiguity rank]\label{def:active-ambiguity}
For $C$ as in \eqref{eq:active-cone-def}, let $V=\spanof(C)$ and put
\[
 J(C)=\bigcup_{\substack{G\in\mathcal G\\|G\cap S(C)|\ge2}}
             (G\cap S(C)).
\]
Define
\begin{equation}\label{eq:active-rho}
 \rho(C)=\dim\spanof\{e_j|_V:j\in J(C)\}.
\end{equation}
\end{definition}

This quantity depends on the cone and the conflict groups, not on the LP
solver's particular basic point.  Universally zero coordinates do not
contribute.  Neither do independent directions carried only by coordinates
outside every active conflict.

It is an exact linear-algebraic quantity once the active-support certificate
has been checked.  If $S=S(C)$, then
\begin{equation}\label{eq:active-dimension}
 V=\{z:Az=0,\ z_j=0\ (j\notin S)\},\qquad
 \dim V=|S|-\rank A_S.
\end{equation}
Indeed, a vector positive on $S$ has a relative neighborhood in
$\ker A_S$ that remains nonnegative, so the cone spans that entire kernel.
Writing $E_J$ for the indicated coordinate rows restricted to $S$ gives
\begin{equation}\label{eq:active-rho-rank}
 \rho(C)=\rank\begin{pmatrix}A_S\\E_J\end{pmatrix}-\rank A_S.
\end{equation}
Rank computation is useful for interpretation and complexity accounting; the
certificate checker does not need to trust a reported rank in order to accept
a positive witness or a negative search tree.

\begin{theorem}[Binary active-face search]\label{thm:active-tree}
Suppose $C$ contains a point of positive objective.  At each nonterminal
positive node, compute complete active support and choose any incompatible
active pair $i,j$.  Recurse on
\[
 C\cap\{x_i=0\}\qquad\text{and}\qquad C\cap\{x_j=0\}.
\]
Terminate at a nonpositive cone or at a positive cone whose active support is
compatible.  This decides admissible positive-objective existence.  Its
complete binary tree has at most $2^{\rho(C)+1}-1$ nodes.  With polynomial
rational LP, its bit cost and certificate length are
$2^{\rho(C)}\poly(L)$, where $L$ is the original encoded input length.
\end{theorem}

\begin{proof}
Every admissible vector satisfies $x_i=0$ or $x_j=0$, so the two children cover
all admissible possibilities.  Their overlap is harmless.  A coordinate
deleted using an active-support dual is zero on the entire current cone;
such a deletion preserves every admissible candidate.  A negative dual
excludes positive objective.  At a compatible positive leaf, the active
primal vector is itself admissible.  These observations prove correctness.

For the bound, let $C'$ be a child whose positive-objective part is nonempty,
and put $V'=\spanof(C')$.  Coordinate restriction cannot create new active
coordinates, so $J(C')\subseteq J(C)$.  Let
\[
 W=\spanof\{e_k|_V:k\in J(C)\}.
\]
In the child $x_i=0$, restriction from $V$ to $V'$ kills $e_i|_V$.  That
functional is nonzero because $i$ was active, and it belongs to $W$ because
$i$ was part of an active conflict.  The restriction map on $W$ therefore
has nontrivial kernel.  Its image has dimension at most $\rho(C)-1$, and
the defining functionals for $\rho(C')$ lie in that image.  Consequently
\[
 \rho(C')\le\rho(C)-1.
\]
The same argument applies to the other child.

If $\rho=0$, there is no active incompatible pair: any active coordinate in
such a pair would be a nonzero defining functional.  Hence every positive
branching node has rank at least one.  Along a path, ranks decrease at each
edge whose child remains positive.  A final negative leaf can occur after
the last rank-one branching node, so the full path still contains at most
$\rho(C)$ branches.  Binary tree counting yields the stated node bound.
Every query and arithmetic witness has polynomial encoded size, and a
negative tree retains both children at each split, giving the time and
certificate estimates.
\end{proof}

The proof makes no assumption that the selected coordinate-zero face drops
dimension by exactly one.  A child may lose several coordinates at once,
which only improves the bound.  It also permits additional early terminal
rules.  In particular, a small ordinary LP may first return a negative dual
or an already admissible basic vector.  Stopping in either case shortens the
tree.  The implementation uses this precheck by default and invokes the
larger active-support LP only for an incompatible positive relaxation.

\subsection{Normal anchors and the precise scope of a conclusion}

For a finite triangulation, choose one zero triangle coordinate at every
global vertex.  Let $m_v$ be the number of local triangle corners representing
vertex $v$.  The complete anchor product has size
\[
 \mathcal A(T)=\prod_v m_v,\qquad \sum_v m_v=4t.
\]
In the one-vertex case, $\mathcal A(T)=4t$.  Restricting the standard matching
matrix and an exact Euler covector to the surviving columns gives the
anchored systems used above.  Every canonical admissible normal vector
belongs to at least one such anchored cone.

If $\rho_{\max}$ bounds the roots that require positive relaxed search, the
complete anchored positive-Euler query therefore has the conditional cost
\begin{equation}\label{eq:active-normal-cost}
 \mathcal A(T)\,2^{\rho_{\max}}\poly(t)
\end{equation}
with a polynomial rational-LP backend.  A negative proof contains every
anchor tuple in the independently reconstructed product, with a complete
binary proof tree for each tuple.  Missing anchors, a missing child, or an
uncertified coordinate deletion invalidate coverage.

A positive conclusion here is a canonical admissible normal vector with
positive Euler characteristic.  A maximal-support vector may have several
components; clearing denominators and dividing by the gcd does not turn it
into an extreme ray or prove connectedness.  The maintained geometric
component/disc endpoint is still required to identify an essential disc.
Conversely, a negative positive-Euler query is one ingredient in a complete
normal-surface recognition architecture, whose provenance and crushing
transitions must also be justified \cite{BurtonOzlen,priorDual}.  The new
module does not convert either status directly into a knot verdict.

\subsection{A linear obstruction already present in genuine root cones}

The active parameter removes identically zero coordinates after restrictions.
At an unrestricted normal root, however, there is a universal source of
activity.  Let $v(T)$ denote the \emph{total} number of global vertices; this
includes boundary vertices and is distinct from the number of interior
vertices used in the batch-selection results.  Fix one zero triangle anchor
per global vertex, and let $u$ be the number of tetrahedra containing at least
one selected anchor.

\begin{theorem}[Unrestricted anchored normal-cone obstruction]
\label{thm:active-root-obstruction}
Allow all three quadrilateral types in every tetrahedron and impose only
the matching, nonnegativity, and chosen triangle-anchor equations.  Then
every quadrilateral coordinate is active, and
\begin{equation}\label{eq:active-root-lower}
 \rho(C)\ge t-u\ge t-v(T).
\end{equation}
This conclusion does not require positive Euler feasibility.  If the
positive-Euler part is nonempty, all quadrilateral coordinates are also
active on that part.  In particular, a one-vertex root satisfies
$\rho(C)\ge t-1$.
\end{theorem}

\begin{proof}
Let $r$ have every triangle coordinate zero and every quadrilateral
coordinate one.  Each triangular face receives one arc of each normal arc
type from these three quadrilaterals.  Thus all face-matching equations
hold.  Every selected triangle anchor vanishes, so $r\in C$ and every
quadrilateral coordinate is active.  Each tetrahedron consequently belongs
to an active three-type conflict.

For a tetrahedron $a$ containing no selected anchor, define $w_a$ to have
value $+1$ in its four triangle coordinates, value $-1$ in its three
quadrilateral coordinates, and zero elsewhere.  A local normal arc count
has the form $T+Q$.  Its change under $w_a$ is therefore $1-1=0$, on every
face of $a$.  This verifies the matching-kernel identity also when faces of
the same tetrahedron are identified.  Both $r$ and $r+w_a$ are nonnegative
and satisfy all chosen anchors.  Hence $w_a\in\spanof(C)$.

The quadrilateral image of $w_a$ is $-1$ precisely in the three coordinates
of tetrahedron $a$.  These images have disjoint supports as $a$ ranges over
the $t-u$ unanchored tetrahedra, so they are linearly independent.  All of
their coordinate functionals occur in the definition of $\rho(C)$ because
all quadrilateral types are active.  Thus the map defining the ambiguity
rank has rank at least $t-u$.  There is at most one selected anchor per
global vertex, and hence $u\le v(T)$, proving the displayed inequality.

Finally, if $p\in C$ has positive Euler characteristic, then
$p+\varepsilon r$ still has positive Euler characteristic for sufficiently
small positive $\varepsilon$, and it has every quadrilateral coordinate
positive.  This proves the last activity assertion.
\end{proof}

The vectors $r$ and $r+w_a$ are incompatible relaxation vectors.  The proof
does not assert that they are embedded surfaces.  Their role is to expose
freedom that an LP relaxation necessarily retains before useful type
restrictions or admissibility implications have been imposed.

The maintained source constructor gives a direct family-specific form of
the obstruction.

\begin{corollary}[The pulling diagram-exterior source]
\label{cor:active-pulling-source}
Let a classical one-component diagram have $n$ crossings and put
$n'=\max\{1,n\}$.  For the maintained pulling construction, with its checked
compact knot-exterior interpretation, every unrestricted anchor root
satisfies
\begin{equation}\label{eq:active-source-linear}
 \rho(C)\ge16n'-2.
\end{equation}
\end{corollary}

\begin{proof}
The source has $4n'$ truncated corner cells and uses five tetrahedra in each
cell, so $t=20n'$.  Each cell contributes two truncation triangles to the
boundary, giving $f_\partial=8n'$.  These counts are part of the maintained
pulling template and source correspondence \cite{ProveItSource}.  The
boundary is a triangulated torus; counting incidences and using its Euler
characteristic gives
\[
 e_\partial=\frac32f_\partial=12n',\qquad
 v_\partial=e_\partial-f_\partial=4n'.
\]
The pulling subdivision introduces no interior vertices beyond the two
finite pole classes of the corner construction.  Its remaining vertices
are boundary cut points.  Thus $v(T)=4n'+2$.  Applying
\eqref{eq:active-root-lower} yields
$\rho(C)\ge20n'-(4n'+2)=16n'-2$.
For $n=0$, the constructor uses its specified one-curl representative of
the crossing-free circle, so the same count with $n'=1$ applies.
\end{proof}

These lower bounds rule out a universal polylogarithmic bound for this
\emph{unrestricted-root} rank along the displayed source family.  They do
not give a lower bound for the recognizer's actual runtime: a negative
ordinary LP or an early admissible point can finish a stage before the
active query runs.  They also do not rule out small parameters after
mandatory propagation, certified type restrictions, or topology changes.
The distinction is concrete.  The earlier continuation proves that its
one-vertex layered Fibonacci family requires no guessed types under the
mandatory rule \cite{priorDual}, whereas
Theorem~\ref{thm:active-root-obstruction} gives a linearly growing
unrestricted active rank on the same family.

\subsection{Implementation contract and research consequence}

The additive implementation has three levels.  The rational-cone primitive
returns a complete active-support pair, an ordinary negative multiplier, or
an inconclusive resource-limit outcome.  The conflict-group search adds
complete binary proof trees.  The normal wrapper reconstructs finite
matching equations, the Euler covector, and the full anchor product, and
binds its final certificate to the source triangulation.  Its independent
checker imports neither the producer nor an LP solver.  Seventeen focused
tests cover arithmetic alternatives, hidden active coordinates, malformed
rationals, incomplete branches, cancellation, anchor/source substitution,
and actual normal positive and negative controls.  In particular, the
positive meridian control is passed onward to the maintained geometric
essential-disc counter.

The mathematical reduction and the executable backend have different roles.
The active-support pair is a short, complete interface for deleting inactive
coordinates.  The binary theorem gives an honest intrinsic bound at any
restricted cone.  The root obstruction identifies an unavoidable linear
source of ambiguity before such restrictions.  Further progress must
therefore explain how geometric information, admissibility implications,
or a controlled transition removes that freedom, or why a verified endpoint
is found before it matters.  Merely replacing allowed masks by active
supports cannot supply the missing global quasi-polynomial theorem.
